\paragraph{Theorem F.2.14 (Proved: four order pairs and eight local spectral pairs).}
Every simply connected two-finite-point datum with $b_2=0$ has
one finite linear order equal to five and the other equal
to four, six, eight or twelve.
After suppressing the common factor $(T-1)^2$, its unordered
pair of finite characteristic polynomials is exactly one
of the following necessary possibilities:
\[
\begin{array}{c|c|l}
 \text{orders}&\text{order-five factor}&\text{other factor}\\\hline
 (5,4)&\Phi_5&\Phi_2^2\Phi_4\\
 (5,6)&\Phi_5&\Phi_2^2\Phi_3,\ \Phi_2^2\Phi_6,\ \Phi_3\Phi_6\\
 (5,8)&\Phi_5&\Phi_8\\
 (5,12)&\Phi_5&\Phi_{12},\ \Phi_3\Phi_4,\ \Phi_4\Phi_6.
\end{array}
 \tag{F.39}
\]
Its rank-five quotient is rationally irreducible.
In particular all order-three pairs, including $(3,4)$,
are impossible.  The table is a necessary spectral filter;
it does not identify the relative rational or integral
positions of the matrices having those polynomials.

\begin{proof}
By F.2.15 the finite actions are locally complex-compatible.
By F.2.9, the rank-six cohomology system $V$ has one
invariant, no coinvariants, and finite fixed ranks two.
Its quotient $W$ has neither invariants nor coinvariants.
Every rational one-dimensional subrepresentation or
quotient character of $W$ is trivial: its two finite
orders divide the coprime affine multiplicities of F.2.7,
so at least one is odd and that character value is one;
the unipotent product forces the other value to be one.
Thus $W$ and $W^*$ have no invariant rational line of any
character.

The space of invariant alternating forms on $W$ has
dimension at most one.  Indeed a nonzero such form cannot
have rank four, because its radical would be a stable
rational line in $W$.  It therefore has rank two and is
a decomposable element of $\bigwedge^2W^*$.
If two independent invariant rank-two forms existed,
their sum would also have rank at most two; hence their
wedge product would be zero.  Their two supporting
planes in $W^*$ would be distinct stable planes meeting
in a line.  That stable rational line contradicts the
preceding paragraph.  This proves the bound without an
irreducibility assumption.

The dual exterior sequence
\[
 0\longrightarrow\bigwedge^2W^*
 \longrightarrow\bigwedge^2V^*
 \longrightarrow W^*\longrightarrow0
\]
and $(W^*)^\Gamma=0$ show that the coinvariant dimension
$b_2'$ of $\bigwedge^2V$ is at most one.
On the other hand its invariant dimension $a_2'$ is zero.
To see this, the corresponding $H^0(R^2g_*\mathbb Q)$
term has no point-supported rational kernel at a finite
fibre or in cusp degree two, and its only possible outgoing
differential has target $H^2(R^1g_*\mathbb Q)=0$.
It therefore survives in total degree two, which vanishes.

Let $d_i$ be the dimension fixed by the $i$th finite
generator on $\bigwedge^2V$.  The cusp fixed dimension is
nine, so the invariant-sheaf Euler identity gives
\[
 -15+d_1+d_2+9
       =a_2'-h^1(j_*\bigwedge^2V)+b_2'\leq1.
 \tag{F.40}
\]
Each finite generator has fixed rank two on $V$.
Its complementary four-dimensional real representation
commutes with a complex structure and has no fixed vectors.
It either has two distinct nonreal inverse pairs, one
nonreal pair with multiplicity two, two negative real
eigenvalues and a nonreal pair, or four negative eigenvalues.
The respective invariant dimensions in degree two are
$3,5,3,7$.  In particular $d_i$ is odd and at least three.
The left side of (F.40) is therefore nonnegative and even,
so it must be zero.  Thus $d_1=d_2=3$.

The complete cyclotomic filter F.2.2 now leaves precisely
the finite polynomials with two fixed eigenvalues and
the indicated degree-two fixed rank three.
Keeping coprime unordered pairs gives the eight rows
expanded in (F.39).  This finite enumeration is reproduced
from the complete saved 38-type list by
\texttt{fourfolds\_order\_pairs.py}.
An order-three generator would instead have characteristic
polynomial $(T-1)^2\Phi_3^2$, hence $d_i=5$, explaining
its exclusion directly.
Every remaining pair contains $(T-1)^2\Phi_5$, so F.2.9
proves rational irreducibility of $W$.
\end{proof}

Theorem F.2.10 supplies an independent spin-representation
proof for three of the excluded order-three pairs and for
the repeated fourth- and sixth-root cases.
Within the third row of the order-six spectral possibilities,
$\Phi_3\Phi_6$, no exclusion follows from F.2.13: that theorem
treats particular rational representations with the first
two order-six spectra.  Likewise F.2.12 excludes one
rational representation of order pair $(5,12)$, not its
entire spectral locus.

